% Book template for the YCB1 exam-guide booklet. Compile with XeTeX (Devanagari support).
% Trade size 6x9 in, as the other books in this workspace. Plain: no heading numbers typed by hand,
% the compiler numbers chapters and sections.
\usepackage[paperwidth=6in, paperheight=9in,
  inner=0.9in, outer=0.7in, top=0.8in, bottom=0.8in]{geometry}
\usepackage{fontspec}
\setmainfont{Segoe UI}
\newfontfamily\devanagarifont[Script=Devanagari]{Nakula}
% Sanskrit or Hindi text: \sa{...}
\newcommand{\sa}[1]{{\devanagarifont #1}}

\usepackage[dvipsnames,table]{xcolor}
\definecolor{ycbsaffron}{RGB}{176, 82, 18}
\definecolor{ycbgrey}{gray}{0.35}

\usepackage{emptypage}
\usepackage{graphicx}
\usepackage{float}
\usepackage{enumitem}
\setlist{nosep}
\usepackage{booktabs}
\usepackage{tcolorbox}
\tcbuselibrary{breakable}

% Running header: chapter title on the left page, section title on the right page.
\usepackage{fancyhdr}
\pagestyle{fancy}
\fancyhf{}
\fancyhead[LE]{\small\textcolor{ycbgrey}{\leftmark}}
\fancyhead[RO]{\small\textcolor{ycbgrey}{\booktitleshort}}
\fancyfoot[C]{\thepage}
\renewcommand{\chaptermark}[1]{\markboth{#1}{#1}}
\renewcommand{\headrulewidth}{0.4pt}
\fancypagestyle{plain}{\fancyhf{}\fancyfoot[C]{\thepage}\renewcommand{\headrulewidth}{0pt}}

\usepackage{titlesec}
\titleformat{\chapter}[display]
  {\normalfont\huge\bfseries\raggedright\color{ycbsaffron}}{\chaptertitlename\ \thechapter}{12pt}{\Huge}
\titlespacing*{\chapter}{0pt}{-10pt}{24pt}
\titleformat{\section}{\normalfont\Large\bfseries\color{ycbsaffron}}{\thesection}{0.6em}{}
\titleformat{\subsection}{\normalfont\large\bfseries}{\thesubsection}{0.6em}{}

\usepackage{xurl}
\usepackage[hidelinks]{hyperref}
\usepackage[nameinlink,capitalise]{cleveref}

% Fewer overfull lines: allow a little extra stretch, and let microtype adjust spacing.
\usepackage[expansion=false]{microtype}
\setlength{\emergencystretch}{3em}
\hbadness=3000

% Figure with its generation prompt below it: \ycbfig{<figure_id>} (ids and prompts come from
% ycb1_book_figures.tex, generated from docs/image_register.csv). When the final picture file does
% not exist yet, an empty frame stands in; the prompt line stays under it either way.
% Order of choice: the final picture (own work, drawn or generated) if its file exists; else the original
% workshop picture, with a "permission to be cleared before print" line when the register marks it high
% risk; else an empty frame. An optional credit line (register column `credit`) is printed when known.
\newcommand{\ycbfigshow}[1]{\includegraphics[width=0.8\linewidth,height=0.3\textheight,keepaspectratio]{#1}}
\newcommand{\ycbfig}[1]{%
  \begin{figure}[H]\centering
  \ifcsname figfile@#1\endcsname\edef\ycbfigfile{\csname figfile@#1\endcsname}\else\def\ycbfigfile{none}\fi
  \ifcsname figorig@#1\endcsname\edef\ycbfigorig{\csname figorig@#1\endcsname}\else\def\ycbfigorig{none}\fi
  \IfFileExists{\ycbfigfile}{\ycbfigshow{\ycbfigfile}\def\ycbfigold{0}}{%
    \IfFileExists{\ycbfigorig}{\ycbfigshow{\ycbfigorig}\def\ycbfigold{1}}{%
      \fbox{\parbox[c][5em][c]{0.8\linewidth}{\centering\footnotesize\textcolor{ycbgrey}{Figure to be made}}}\def\ycbfigold{0}}}\par\smallskip
  {\footnotesize\itshape\textcolor{ycbgrey}{%
    \ifcsname figcredit@#1\endcsname Credit: \csname figcredit@#1\endcsname.\ \fi
    \if1\ycbfigold\ifcsname figclear@#1\endcsname Original workshop picture; permission to be cleared before print.\ \fi\fi
    \ifcsname figdrawn@#1\endcsname Line diagram drawn for this book.\else
    Prompt: \ifcsname figprompt@#1\endcsname\csname figprompt@#1\endcsname\else to be written (see the image register, #1)\fi\fi}\par}
  \end{figure}}

% Draft label for a chapter written from a notes block whose Fact check is still pending
% (author's decision 2026-11-24: write it now, labelled; remove the box when the block is checked).
\newtcolorbox{unverified}{breakable, colback=yellow!8, colframe=orange!70!black, title=Not yet fact-checked, fonttitle=\bfseries}

% Boxes used by the chapters: a cautions box and a quick-recall box.
\newtcolorbox{caution}{breakable, colback=red!4, colframe=red!55!black, title=Caution, fonttitle=\bfseries}
\newtcolorbox{recall}{breakable, colback=ycbsaffron!6, colframe=ycbsaffron, title=Quick recall, fonttitle=\bfseries}
